\chapter{Wang Wei: Landscape and Emptiness}

\section{Life}

Wang Wei \zht{王維} (701?--761) was the model High Tang gentleman: a \emph{jinshi}
graduate at about twenty, a court official for most of his life (ending as Assistant
Director of the Right, so he is often called \zht{王右丞}), a celebrated musician,
calligrapher and painter, and a devout lay Buddhist.

\begin{itemize}
  \item His courtesy name Mojie \zht{摩詰} together with his given name Wei \zht{維} makes
        \zht{維摩詰}, Vimalakīrti, the wise layman of the Buddhist \emph{Vimalakīrti
        Sutra}: a sign of the Buddhist identity he cultivated.
  \item After his wife's death (c.\ 730) he never remarried, and he lived increasingly like a
        monk.
  \item He owned an estate at Wangchuan \zht{輞川} in the Zhongnan mountains south of
        Chang'an; with his friend Pei Di he wrote a sequence of twenty quatrains on its
        sites, the \emph{Wangchuan Collection} \zht{輞川集}.
  \item In 756 the rebel An Lushan captured Chang'an; Wang Wei was forced to serve in the
        rebel government. After the recapture he was nearly punished but was pardoned,
        partly thanks to a poem of loyal grief he had written in captivity and partly
        thanks to his brother's intercession.
\end{itemize}

\begin{important}[``Poems in paintings'']
Su Shi (Song dynasty) famously wrote: \zht{味摩詰之詩，詩中有畫；觀摩詰之畫，畫中有詩。}
``Savour Mojie's poems and there are paintings in them; look at Mojie's paintings and there
are poems in them.'' Wang Wei was later revered as the founder of the ``Southern school''
of literati monochrome landscape painting. His poems are visual, spatial and quiet.
\end{important}

\section{Key ideas for reading Wang Wei}

\begin{center}
\begin{tabular}{lp{0.65\textwidth}}
\toprule
Concept & In Wang Wei \\
\midrule
\zht{空} \emph{kōng} ``empty'' & His signature word: the empty mountain. In Buddhism, \emph{śūnyatā}:
  all things lack fixed self-nature. The landscape is both literally deserted and
  metaphysically empty. \\
\zht{靜} \emph{jìng} ``still'' & Stillness revealed by a small sound or movement. \\
Absence of self & Few first-person pronouns; the observer dissolves into what is
  observed. \\
Light & Slanting evening light, moonlight: the changing play of appearance. \\
Understatement & No explicit doctrine; the Buddhist meaning is in the perception itself. \\
\bottomrule
\end{tabular}
\end{center}

\section{``Deer Park'' \zht{鹿柴}}

From the \emph{Wangchuan Collection}. (\zht{柴} is read \emph{zhài} ``fence, enclosure''
here.)

\begin{center}
\begin{tabular}{p{0.22\textwidth}p{0.28\textwidth}p{0.4\textwidth}}
\toprule
Original & Pinyin & Gloss \\
\midrule
\zht{空山不見人，} & kōng shān bú jiàn rén & empty mountain not see person \\
\zht{但聞人語響。} & dàn wén rén yǔ xiǎng & only hear person speech echo \\
\zht{返景入深林，} & fǎn yǐng rù shēn lín & returning light enter deep forest \\
\zht{復照青苔上。} & fù zhào qīng tái shàng & again shine green moss on \\
\bottomrule
\end{tabular}
\end{center}

My translation:

\begin{example}[``Deer Park'']
Empty mountain: no one to be seen, \\
only the echo of someone's voice. \\
Light, turning back, comes deep into the wood \\
and shines once more upon the green moss.
\end{example}

\subsection{Analysis}
\begin{itemize}
  \item \textbf{Senses.} Lines 1--2: sight fails (no one seen), hearing remains (a voice
        echoes). Lines 3--4: sight returns, but only a patch of light on moss.
  \item \textbf{\zht{返景}} (\zht{景} = \zht{影} \emph{yǐng}, ``light, sunlight''): the
        sunset light that slants back in under the canopy. \zht{復} ``again'' implies the
        light was here before (morning) and has returned; time is cyclical.
  \item \textbf{Emptiness.} The mountain is ``empty'' though people are nearby: the voices
        only make the emptiness felt. The moss lit for a moment is an image of
        transient appearance.
  \item \textbf{Form.} Five-character quatrain, rhymes \emph{xiǎng / shàng} (oblique
        rhyme), so not strictly regulated: an ``old-style'' quatrain.
  \item \textbf{Translation problem.} No subject: who hears? ``I'' makes it a hiker's
        anecdote; leaving it out keeps the impersonal perception. The essay collection
        \emph{Nineteen Ways of Looking at Wang Wei} (Eliot Weinberger, 1987) compares many
        English versions of this one poem; it was on reserve for the close-reading paper.
\end{itemize}

\section{``Bird-Call Valley'' \zht{鳥鳴澗}}

\begin{center}
\begin{tabular}{p{0.3\textwidth}p{0.6\textwidth}}
\toprule
\zht{人閒桂花落，} & People at rest; osmanthus blossoms fall. \\
\zht{夜靜春山空。} & The night is still; the spring mountain empty. \\
\zht{月出驚山鳥，} & The moon comes out and startles the mountain birds: \\
\zht{時鳴春澗中。} & now and then they call in the spring ravine. \\
\bottomrule
\end{tabular}
\end{center}

\textbf{Analysis.} So quiet that one can sense blossoms falling. The \emph{moon rising}
is enough to startle the birds: the stillness is broken by light, then by sound, and the
cry deepens the silence (a principle also of classical Japanese poetry influenced by
Wang Wei). \zht{閒} ``at leisure'' describes both people and mind.

\section{``Lodge in the Bamboo Grove'' \zht{竹里館}}

\begin{center}
\begin{tabular}{p{0.3\textwidth}p{0.6\textwidth}}
\toprule
\zht{獨坐幽篁裏，} & Sitting alone in the hidden bamboo, \\
\zht{彈琴復長嘯。} & I play the zither, then whistle long. \\
\zht{深林人不知，} & Deep in the wood, no one knows; \\
\zht{明月來相照。} & the bright moon comes to shine on me. \\
\bottomrule
\end{tabular}
\end{center}

\textbf{Analysis.} The ``long whistle'' (\zht{嘯}) was a Daoist breathing practice
associated with eccentric recluses. ``No one knows'' (\zht{人不知}) echoes the Confucian
complaint of the unrecognised scholar, but here it is turned around: being unknown is
peace, and the moon is companion enough. Compare Li Bai's drinking with the moon
(next chapter), which is far less serene.

\section{``Autumn Evening in the Mountains'' \zht{山居秋暝} (regulated octave)}

\begin{center}
\begin{tabular}{p{0.3\textwidth}p{0.6\textwidth}}
\toprule
\zht{空山新雨後，} & In the empty mountains, after fresh rain, \\
\zht{天氣晚來秋。} & the air at evening feels like autumn. \\
\zht{明月松間照，} & The bright moon shines among the pines; \\
\zht{清泉石上流。} & the clear spring flows over the stones. \\
\zht{竹喧歸浣女，} & Bamboo rustles: washing-girls going home; \\
\zht{蓮動下漁舟。} & lotus stirs: a fishing boat going down. \\
\zht{隨意春芳歇，} & Let the spring flowers fade as they will: \\
\zht{王孫自可留。} & the young lord can still stay here. \\
\bottomrule
\end{tabular}
\end{center}

\subsection{The parallel couplets}

\begin{center}
\begin{tikzpicture}[ch/.style={draw, minimum size=0.75cm},
  gl/.style={font=\scriptsize}, cat/.style={font=\scriptsize, text=blue!60!black}]
  \foreach \c [count=\i] in {明,月,松,間,照} { \node[ch, fill=yellow!20] (a\i) at (1.2*\i,3.4) {\zht{\c}}; }
  \foreach \c [count=\i] in {清,泉,石,上,流} { \node[ch, fill=cyan!15] (b\i) at (1.2*\i,2.2) {\zht{\c}}; }
  \foreach \i in {1,...,5} { \draw[gray, <->] (a\i) -- (b\i); }
  \foreach \t [count=\i] in {adj.,noun,noun,locative,verb} { \node[cat] at (1.2*\i,4.15) {\t}; }
  \foreach \c [count=\i] in {竹,喧,歸,浣,女} { \node[ch, fill=green!15] (c\i) at (1.2*\i,0.9) {\zht{\c}}; }
  \foreach \c [count=\i] in {蓮,動,下,漁,舟} { \node[ch, fill=pink!25] (d\i) at (1.2*\i,-0.3) {\zht{\c}}; }
  \foreach \i in {1,...,5} { \draw[gray, <->] (c\i) -- (d\i); }
  \foreach \t [count=\i] in {plant,verb,verb,activity,noun} { \node[cat] at (1.2*\i,-1.05) {\t}; }
  \node[gl, anchor=west, align=left] at (6.8,2.8) {sky / earth\\ light / water\\ still / moving};
  \node[gl, anchor=west, align=left] at (6.8,0.3) {effect before cause:\\ we hear the bamboo\\ before we see the girls};
\end{tikzpicture}
\end{center}

\subsection{Analysis}
\begin{itemize}
  \item \textbf{Couplet 2} (lines 3--4) is a perfectly balanced landscape: above and below,
        moonlight and water, pine and stone. Nothing human.
  \item \textbf{Couplet 3} (lines 5--6) brings people in, but \emph{indirectly}: first the
        rustling and the stirring lotus, only then the girls and the boat. The syntax
        follows the order of perception: a famous example of Wang Wei's ``inverted''
        lines.
  \item \textbf{Ending.} An allusion reversed: the \emph{Chuci} poem ``Summoning the
        Recluse'' calls the ``young lord'' (\zht{王孫}) to come \emph{back} from the
        mountains, which are dangerous. Wang Wei says: stay; the mountain is home.
  \item \textbf{Form.} Pentasyllabic regulated octave; rhymes \emph{qiū, liú, zhōu, liú}
        (level).
\end{itemize}

\section{A farewell poem}

``Seeing Off Yuan the Second on a Mission to Anxi'' \zht{送元二使安西} was so popular it
was sung at farewells as the ``Song of the Three Refrains at Yang Pass'' (\zht{陽關三疊}):

\begin{center}
\begin{tabular}{p{0.38\textwidth}p{0.52\textwidth}}
\toprule
\zht{渭城朝雨浥輕塵，} & In Weicheng the morning rain dampens the light dust; \\
\zht{客舍青青柳色新。} & by the inn, green, green, the willows are fresh. \\
\zht{勸君更盡一杯酒，} & I urge you, sir, to drain one more cup of wine: \\
\zht{西出陽關無故人。} & west of Yang Pass you will have no old friends. \\
\bottomrule
\end{tabular}
\end{center}

\begin{remark}
Every element is conventional for a farewell (rain, willows, wine, the frontier pass),
yet the poem is completely fresh. The willow (\zht{柳} \emph{liǔ}, punning on \zht{留}
\emph{liú} ``stay'') is the hidden plea; the open statement is just ``one more cup''.
Yang Pass was the gateway to the western regions (modern Gansu), the edge of the Chinese
world.
\end{remark}

\section*{Summary}
\begin{itemize}
  \item Wang Wei: High Tang official, painter, musician, lay Buddhist (Mojie =
        Vimalakīrti); Wangchuan estate; compromised by the An Lushan rebellion.
  \item Signature words \zht{空} empty, \zht{靜} still; absence of self; perception
        presented in sequence.
  \item ``Deer Park'': voice without person, returning light on moss.
  \item ``Autumn Evening in the Mountains'': textbook parallel couplets; reversal of the
        ``Summoning the Recluse'' allusion.
  \item ``Poems in paintings, paintings in poems'' (Su Shi).
\end{itemize}
